\section{Related Work}

\subsection{The Governance-Level Positioning This Paper Inherits}

SIDS 2026 \cite{katrakazas2026} situates BENEFITS against three bodies of prior work. The first is Data Space governance architecture: the IDSA Reference Architecture Model \cite{idsram4} establishes connector-level usage control as the primary sovereignty mechanism for Data Spaces, and the Gaia-X Architecture Document \cite{gaiax} extends this with certification-based obligations around portability, interoperability, and transparency. SIDS 2026 \cite{katrakazas2026} reads both as designed for organizational actors with legal standing and negotiating capacity, leaving individual data subjects, particularly those in precarious situations, as governance beneficiaries rather than governance participants. It reads the EU Data Governance Act and Data Act the same way: as formalizing intermediary roles and access rules at the level of institutions, not individuals.

The second body of prior work is citizen data sovereignty projects, anchored on DECODE \cite{decode2019}, a blockchain-based architecture with attribute-based cryptography for privacy-preserving identity management, piloted in Barcelona and Amsterdam. SIDS 2026 \cite{katrakazas2026} credits DECODE with demonstrating that decentralized architectures can give users granular control over personal data, while noting that its always-on network requirement and its assumption of digital literacy and device ownership make it unsuitable for the offline-first contexts the BENEFITS participant population needs. The third body is participatory sensing and community data governance: SPPEAR \cite{gisdakis2014}, community-based participatory research methods \cite{vaughn2020}, equitable participatory design (Harrington et al., 2019; Giorgi et al., 2022), and civic-tech platforms such as Open Humans \cite{openhumans2019} and the MyData operator model \cite{mydata2020}. SIDS 2026 \cite{katrakazas2026} reads these as treating governance as external to the architecture, enforced by the system operator rather than verifiable by participants, and reads recent informed-consent-design work \cite{mcinnis2024} as still treating consent as a communication and comprehension problem rather than a verifiability problem.

SIDS 2026's own comparison table \cite{katrakazas2026} carries this positioning into a direct feature comparison against IDSA RAM, Gaia-X, DECODE, SPPEAR, Open Humans, MyData, Chereja et al.'s by-design privacy-risk architecture \cite{chereja2025}, and McInnis et al.'s consent-UX work, scored across sovereignty enforcement locus, target population, offline capability, community governance, and privacy mechanism. Its prose conclusion is that BENEFITS is distinguished by the combination of sovereignty enforcement at the participant device, offline-first multi-channel collection, community-embedded governance with technical veto power, and explicit design for high-distrust populations, and that no existing framework in the table addresses this intersection. The table itself carries a fifth dimension beyond that prose summary, though it is a narrower distinction than a first read suggests: three comparators, DECODE (attribute-based cryptography), SPPEAR (pseudonyms), and Chereja et al. (a risk index combining a privacy risk expansion factor and a privacy exposure index), also name a concrete privacy technique, not merely a policy, certification, or UX category. What distinguishes BENEFITS is not being the only row with a concrete technique, but combining two concrete techniques, k-anonymity and differential privacy, where each of those three comparators names only one. Taken together, the table's own logic is that BENEFITS is the sole entry combining edge enforcement, vulnerable-population targeting, offline support, community steward veto, and the pairing of k-anonymity with DP.

\subsection{From Governance Comparators to Technical Peers}

That comparison operates entirely at the level of architecture and governance model. Every row in SIDS 2026's table, including the BENEFITS row itself, is scored on what a design proposes to do, not on what an implementation has been shown to do. This was appropriate to what SIDS 2026 was: its own "Technical Implementation Status" places the steward-quorum release gate and the differential-privacy pipeline in the design and specification phase, pending pilot deployment. A framework that exists as specification cannot honestly be evaluated against implemented systems, so a comparison against other governance proposals was the correct scope for that paper.

This paper's contribution changes that scope. The BENEFITS reference architecture is now a running system: an EDC-based connector carrying policy, quorum, and privacy-pipeline extensions; a hash-chained audit log; a consent layer synchronizing PouchDB to CouchDB with edge-level enforcement; and six mechanism-level experiments measuring what each of these pieces actually costs and where it breaks. Positioning that artifact against governance frameworks alone would leave its real technical neighbors unexamined. The four clusters below extend SIDS 2026's related-work section one layer down, from architectures that propose sovereignty properties to systems that implement and, in some cases, measure the same class of property: Data Space connector implementations with a reported technical evaluation, Merkle-tree and hash-chain tamper-evident logging systems, differential-privacy and k-anonymity systems with an explicit budget accountant, and threshold-cryptography or quorum-authorization systems for data governance.

\subsection{Data Space Connector Implementations}

The closest technical neighbor to BENEFITS's own EDC deployment is a small and largely non-benchmark literature. \cite{dam2023} survey open-source IDS-compliant dataspace connector implementations, comparing their architecture, data model, and usage-control language support against the IDS Reference Architecture Model's requirements. The survey is qualitative rather than empirical, but it establishes the same premise this chapter extends: which connectors are real, running software with verified enforcement, as opposed to software that merely claims IDS compliance, is itself an open comparative question in the connector literature.

\cite{mugge2023} go further, implementing a Catena-X connector deployment (built on Eclipse Dataspace Components, the same connector family BENEFITS uses) for cross-company end-of-life-vehicle decision support, and validating it against an actual cross-company data flow between named industrial roles. It is the strongest match in this cluster for "a real EDC-based deployment with a reported technical evaluation," though its evaluation is a use-case validation rather than a latency or throughput benchmark. \cite{kalogeropoulos2023}, in an arXiv preprint without a confirmed peer-reviewed venue, report a working open-source connector-as-a-service prototype integrated into a Multi-access Edge Computing architecture, with experiments measuring functionality and scalability of connector-mediated data exchange. It is one of the few sources in this cluster with any quantitative evaluation component, even though that evaluation sits at the edge-computing layer rather than at the connector's own enforcement layer.

\cite{norikane2025}, published in this paper's own target venue, model a data-space-based manufacturing supply chain using the Unified Architecture Framework and run scenario-based simulations comparing conventional manufacturing recovery against a data-space-coordinated alternative under physical-security and cybersecurity disruption. Their finding, that the data-space-coordinated alternative outperforms conventional recovery on production volume and cost under simulated disruption, evaluates supply-chain resilience outcomes rather than the connector's own technical properties, making it a weaker match for this cluster's core question than Mügge et al. or Kalogeropoulos et al. It is included because it is verified, technically evaluated prior work on a built data-space deployment appearing in MDPI Systems itself.

None of these four sources reports a rigorous connector-level latency or throughput benchmark under load. This is worth stating plainly rather than implying a body of comparator work exists where it does not: EDC performance evaluation is thin in the literature generally, and we are aware of no connector performance study in the paper's own target venue either. Given how quickly EDC-based deployments are spreading (Catena-X, Manufacturing-X, Mobility Data Space among them), this is a real and citable gap that remains open in this manuscript as well: none of the six mechanism-level experiments introduced in Chapter 4 benchmarks the connector's own request-handling latency or throughput, since each targets a different boundary of the architecture, the device-side consent filter, the audit chain, the privacy pipeline, or the steward-quorum policy, rather than the connector's own performance under load.

\subsection{Merkle-Tree and Hash-Chain Tamper-Evident Logging}

The audit log SIDS 2026 describes as translating formal sovereignty claims into observable, verifiable controls rests on a mechanism the base paper never traces to a source. \cite{schneier1999} supply that missing citation: their scheme for secure audit logs combines chained one-way hashes with forward-secure message authentication codes evolved under a one-way key-update function, so that log entries written before a compromise are unreadable to an attacker who later controls the logging machine, a property of the forward-secure key evolution, and undetectably unmodifiable, a property of the hash chaining, because altering any already-hashed entry breaks the chain from that point forward. This retroactive-tamper-detection guarantee is exactly the property a Merkle-epoch audit chain generalizes.

\cite{crosby2009} improve on chained-hash schemes with a Merkle-tree-based, history-tree-style structure producing logarithmic-size membership and consistency proofs rather than linear-size ones, and they evaluate proof size and computation cost empirically as log size scales, establishing proof-size-versus-log-size as the metric of interest for this entire class of system. Certificate Transparency \cite{laurie2013}, originally specified as RFC 6962 and later superseded by RFC 9162 (Certificate Transparency Version 2.0), is the most widely operational instance of this design pattern: it defines standardized Merkle audit paths and consistency proofs for publicly logging TLS certificate issuance, and has been deployed at internet scale since publication; Chrome has enforced compliance since 2018 and Firefox since 2025, so the major browsers now require it for publicly trusted TLS certificates. CONIKS \cite{melara2015} applies the same family of structure to key transparency, letting end users audit their own key bindings without a global monitor, and reports measured bandwidth cost directly, under 20 kB per user per day even for billion-user providers, and roughly 2.5 kB per provider per day for collective non-equivocation auditing. \cite{hu2021} push the same evaluation shape further with Merkle², a forest of chronological Merkle trees nesting prefix-trees of child logs that decouples update latency from tree depth, reporting update propagation in roughly one second and support for on the order of 100 times more users than prior transparency logs at comparable auditor cost.

These five sources form a coherent methodological lineage, from the original hash-chain guarantee through tree-structured proofs to measured proof-size and latency behavior at scale, and they collectively anchor this paper's own audit-chain design and its proof-size and inclusion-verification-time measurements. But among the sources in this cluster that report an explicit empirical scale sweep, Crosby and Wallach's sweep of log size and CONIKS's and Merkle²'s sweeps of user population, the variable measured against system scale is log size or user population, not epoch size. This paper's own audit-chain experiment sweeps epoch size specifically, batching entries into periodically rooted Merkle epochs and measuring how proof size and verification time grow as epoch size changes. None of those three sources sweeps epoch size directly, which is a genuine point of methodological departure rather than direct precedent.

\subsection{Differential Privacy and k-Anonymity with Explicit Budget Accounting}

BENEFITS combines a structural k-anonymity stage with a Laplace-noise differential-privacy mechanism governed by a per-site budget accountant. \cite{abadi2016} supply the canonical worked example of what a budget-accounting mechanism paired with an explicit accuracy measurement looks like: their moments accountant tracks cumulative privacy loss across many training steps of DP-SGD, and they report the resulting accuracy cost directly, 97 percent on MNIST and 73 percent on CIFAR-10 under specific epsilon and delta budgets. \cite{mironov2017} supplies the theoretical machinery much of the modern budget-accounting literature, including systems structurally similar to a fixed per-site total budget, is built on: Rényi differential privacy, a relaxation whose composition rule is substantially tighter than naive composition and that recovers and generalizes the moments-accountant technique Abadi et al. use above, so that accountant is properly understood as a special case of Rényi differential privacy accounting rather than a separate approach RDP simply supersedes.

\cite{hay2016} formalize the empirical side of this question with DPBench, an open benchmark establishing evaluation principles for comparing DP algorithms' accuracy under varying privacy budgets and query workloads over range-query tasks, producing concrete accuracy-versus-epsilon curves across multiple mechanisms under one shared methodology. It remains the closest methodological precedent for reporting an accuracy-versus-epsilon curve, though at range-query workloads rather than the small-population social-service aggregates this paper's own privacy-pipeline experiment targets. \cite{rogers2021} describe LinkedIn's Audience Engagements API, a production, deployed budget-management service enforcing a strict differential-privacy budget on returned query results at industrial scale, with explicit engineering tradeoffs for keeping results useful within that fixed budget. It is the strongest evidence in this cluster that a budget accountant with a measured utility cost is a real, deployed design pattern and not only a theoretical construction, and it is structurally the closest analogue to BENEFITS's own per-site budget accountant. \cite{abowd2022} document the largest deployed application of a formal privacy-budget-accounting framework at national scale: the 2020 Census Disclosure Avoidance System's TopDown algorithm, built on zero-concentrated differential privacy, including the policy process for allocating privacy-loss budget across query tables and the resulting accuracy and bias tradeoffs against the pre-2020 disclosure-avoidance approach. \cite{prasser2020} describe ARX, an actively maintained, open-source tool implementing k-anonymity and related models (l-diversity, t-closeness, k-map) alongside a differential-privacy mode, with built-in utility metrics for the information-loss cost of a chosen anonymization configuration, the closest existing tool that combines structural anonymization and formal DP the way BENEFITS's own privacy pipeline does.

Read together, this cluster splits cleanly between algorithmic and theoretical accounting work (Abadi et al., Mironov) and production-deployment reports at very large scale (Rogers et al., Abowd et al.), with comparatively few mid-scale, small-population empirical evaluations in between. SIDS 2026 \cite{katrakazas2026} itself raises this as an open governance concern, noting that differential privacy can render aggregate statistics analytically meaningless at the low population volumes typical of individual service sites. To our knowledge, no peer-reviewed source closes that concern; DPBench is the nearest methodological precedent, but at a different workload scale entirely. This paper's own budget-versus-utility experiment sits in that empirical gap directly.

\subsection{Threshold Cryptography and Quorum Authorization for Data Governance}

The community-steward quorum that gates data release in BENEFITS has its closest technical analogue in threshold-cryptography systems for data governance, but this is the thinnest of the four clusters surveyed for this chapter, with only three sources meeting inclusion criteria. \cite{kokoriskogias2020} describe CALYPSO, a system for "on-chain secrets," confidential data whose release requires a threshold of a trustee committee, with an atomic, auditable access trace and user-controlled consent management. CALYPSO benchmarks transaction latency directly against committee size, reporting latency scaling from roughly 0.2 to 8 seconds as the trustee committee grows from 16 to 128 members, an explicit, measured tradeoff between quorum size and system latency and availability. It is the single strongest match in this cluster and the direct comparator for this paper's own quorum-availability experiment, which measures the availability cost of enforcing the veto guarantee SIDS 2026's own specification calls for against the built system's veto-blind policy.

\cite{raisaro2019} describe MedCo, a deployed system in which decrypting pooled clinical and genomic query results requires the collective participation of all federated clinical sites holding key shares, so that no single site can unilaterally release data. This is structurally close to a steward-quorum release gate in an institutional rather than community setting, though MedCo's published evaluation centers on latency and scalability as the number of participating sites grows rather than on a quorum-size-versus-failure-rate curve, and its decryption gate requires universal rather than threshold participation, a weaker match on the specific availability-tradeoff dimension than CALYPSO. \cite{zhao2024} propose a threshold traceable delegation-authorization scheme for electronic health record sharing, in which multiple authorized delegates must jointly complete a threshold signature operation before an out-of-policy access request is granted, with traceability of which delegates participated, and report experiments on the scheme's practicability. Whether those experiments include an explicit availability-versus-threshold-size tradeoff curve, as opposed to latency or overhead benchmarks alone, is unclear from the published description; it is cited here only for the mechanism it proposes, not for a specific quantitative claim.

Most threshold-cryptography papers report latency or throughput benchmarks under varying threshold sizes rather than a quorum-size-versus-failure-rate curve of the kind this paper's own quorum experiment produces. Only CALYPSO clears that bar cleanly. That leaves this paper's quorum-availability result among the more direct empirical availability-versus-security tradeoff measurements in this specific application area, community and steward-based data governance, rather than simply in the broader threshold-cryptography literature, where the three sources cited here should be read as an honest reflection of how little empirical prior art exists at this intersection, not as an exhaustive survey of the field.

\subsection{Updated Positioning: Verified, Not Just Architected}

One further source sits outside all four clusters but belongs in this positioning nonetheless. \cite{albalwy2026}, also published in this paper's target venue, describes ZK-EHR, a decentralized access-control framework for cross-institutional electronic health record sharing that integrates zk-SNARK zero-knowledge proofs with blockchain smart contracts to verify policy compliance without revealing a requester's role, affiliation, or credentials. It uses neither threshold cryptography nor a Merkle audit log nor a differential-privacy budget, so it does not fit the Merkle-audit-log cluster of Section 2.4, the privacy-budget cluster of Section 2.5, or the threshold-cryptography cluster of Section 2.6 on mechanism, but it makes the same structural move this paper makes: replacing an asserted governance property, in Albalwy's case policy compliance, with one that is cryptographically checkable, in a health-data-sharing context comparable to BENEFITS's own vulnerable-population setting.

Across all four clusters and this additional source, no single system combines the three mechanisms this paper measures together. CALYPSO measures quorum-size tradeoffs but has no audit chain or DP budget; Merkle² and CONIKS measure audit-log proof-size and latency but have no consent-quorum or privacy-budget component; Rogers et al. and Abadi et al. measure budget-versus-utility tradeoffs but have no quorum or tamper-evident log. Each cluster's strongest source addresses one mechanism in isolation. This is itself evidence for this paper's own contribution: no source surveyed here evaluates an audit chain, a privacy-budget accountant, and a quorum-based release gate together, under one built system, the way this paper's six mechanism-level experiments do.

SIDS 2026's comparison table \cite{katrakazas2026} positioned BENEFITS against governance frameworks along five dimensions: sovereignty enforcement locus, target population, offline capability, community governance, and privacy mechanism, concluding that BENEFITS was the sole entry combining edge enforcement, vulnerable-population targeting, offline support, community steward veto, and k-anonymity plus DP. Every comparator in that table, and every entry in it including BENEFITS itself at the time it was written, was scored on architectural design rather than measured behavior. The technical neighbors surveyed in this chapter change what a sixth dimension to that comparison would need to ask: not whether a system claims a sovereignty property, but whether that property has been measured under stress and reported honestly, including where it costs more than expected or degrades. CALYPSO reports that quorum size trades directly against latency. Rogers et al. report explicit engineering tradeoffs for preserving utility within a fixed, enforced production budget, and Abowd et al. report a real, publicly contested accuracy cost at national statistical-agency scale, documented and defended against public and academic criticism. Crosby and Wallach and their successors report that audit-log proof size grows predictably with what is being proven. None of these findings comes from a system that also implements community-steward governance for a vulnerable population, offline-first participation, or the specific combination SIDS 2026's own table established as distinguishing BENEFITS from every governance-level comparator. This paper's contribution is to hold both properties at once: a system with the same governance-level distinguishing combination SIDS 2026 already established, now evaluated with the same empirical rigor this chapter's technical clusters bring to narrower single-mechanism systems. Verified, not just architected, is the dimension this paper adds to SIDS 2026's own positioning.
